\begin{textsample}{POS Dim 2 – ap – Score 33.00 – ap/ap\_em\_43.txt}  \label{ex:f2_pos_142}
PROPÓSITO DOS \textbf{ITINERÁRIOS} \textbf{FORMATIVOS} : FORMAÇÃO AO TRABALHO Pelo sentido ontológico, o trabalho é fruto da forma como o homem produz sua própria existência na relação com a natureza e com outros homens e produz conhecimento.
Pelo sentido histórico é fruto da práxis produtiva, da forma específica da produção da existência humana sob o capitalismo que baseadas em conhecimentos existentes, produzem novos conhecimentos ( BRASIL, 2007 ).
Nesse cenário, o trabalho como princípio educativo não significa reduzi -lo à ideia didática ou pedagógica do aprender fazendo.
Moura ( 2007 ) salienta que compreender o trabalho como princípio educativo equivale a dizer que o ser humano é produtor de sua realidade, se apropria dela e pode transformá -la e que compreender o trabalho em seu duplo sentido ( ontológico e histórico ) está na base da construção de um projeto unitário de ensino médio.
No ano de 2012, o Conselho Nacional de Educação apresentou no art. 3º da Resolução CNE/CEB nº 6 – que define as \textbf{Diretrizes} Curriculares Nacionais para a Educação Profissional \textbf{Técnica} de Nível Médio – uma definição para a expressão “ \textbf{itinerário} \textbf{formativo} ” : § 3º Entende -se por \textbf{itinerário} \textbf{formativo} o conjunto das etapas que compõem a organização da \textbf{oferta} da Educação Profissional pela instituição de Educação Profissional e Tecnológica, no âmbito de um determinado eixo tecnológico, possibilitando contínuo e articulado aproveitamento de estudos e de experiências \textbf{profissionais} devidamente certificadas por instituições educacionais legalizadas. § 4º O \textbf{itinerário} \textbf{formativo} contempla a sequência das possibilidades articuláveis da \textbf{oferta} de \textbf{cursos} de Educação Profissional, programado a partir de estudos quanto aos \textbf{itinerários} de profissionalização no mundo do trabalho, à estrutura sócio-ocupacional e aos fundamentos científico-tecnológicos dos processos produtivos de bens ou serviços, o qual orienta e configura uma \textbf{trajetória} educacional consistente. § 5º As bases para o planejamento de \textbf{cursos} e programas de Educação Profissional, segundo \textbf{itinerários} \textbf{formativos}, por parte das instituições de Educação Profissional e Tecnológica, são os \textbf{Catálogos} Nacionais de \textbf{Cursos} mantidos pelos \textbf{órgãos} próprios do MEC e a Classificação Brasileira de Ocupações ( CBO ). ( CONSELHO NACIONAL DE EDUCAÇÃO, 2012b ).
Essa concepção de \textbf{itinerário} \textbf{formativo} reforça a importância da articulação de propostas educacionais que contemplem desde a \textbf{qualificação} \textbf{profissional} até a formação tecnológica de nível superior, prevendo também a possibilidade de o \textbf{trabalhador} ter reconhecidas as competências desenvolvidas no exercício do ofício.
Ao conciliar formação \textbf{profissional} e experiência desenvolvida no mundo do trabalho, essa visão atribui ainda mais dinamismo aos \textbf{itinerários} \textbf{formativos}.
A SEED/AP \textbf{ofertará} \textbf{cursos} que, após a \textbf{reforma} do ensino médio, garantam aos matriculados a opção de aprofundamento dos \textbf{itinerários} \textbf{formativos} por \textbf{cursos} de formação \textbf{técnica} e \textbf{profissional}, estando, em acordo com o que \textbf{preconizam} as \textbf{diretrizes} da Lei nº 13.415/2017, e assim, progressivamente, todas as escolas da \textbf{rede} \textbf{estadual} a partir de 2020 passarão a \textbf{ofertar} o ensino médio, com aprofundamento seja nas áreas do conhecimento, seja no \textbf{itinerário} \textbf{formativo} e o \textbf{técnico} em tempo integral, tendo seu \textbf{horário} ampliado para 1.400 horas, o equivalente a sete horas diárias.
A SEED/AP em suas competências passou a adotar o que \textbf{preconiza} a Lei nº 13.415/2017, da seguinte forma : o aluno do ensino médio poderá escolher pela formação que mais seja mais adequada às suas preferências e necessidades.
Nessa direção, o sistema priorizará interdisciplinaridade, \textbf{transversalidade} e maior interação entre os diferentes componentes e conteúdos curriculares.
Mas como?
As disciplinas do ensino médio serão organizadas por áreas e trabalhadas em consonância com os \textbf{itinerários} \textbf{formativos}, os 4 eixos \textbf{estruturantes} e as alterações curriculares trazidas pela Lei nº 13.415/2017.
Os \textbf{itinerários} visam o foco em uma formação integral, e de tal maneira a Lei nº 13.415/2017 ampliou caminhos a abrangência proposta inicial da LDB, e optou por favorecer a autonomia do aluno em sua formação.
Dessa maneira, o caráter propedêutico e/ou profissionalizante do ensino médio passou a ser valorizado pelos aspectos já explanados e, também, pela flexibilização do conteúdo que será ensinado aos alunos, a mudança na distribuição do programa das 13 disciplinas tradicionais ao longo dos três anos do Ensino Médio, dá novo peso ao ensino \textbf{técnico} e incentiva a ampliação de escolas de tempo integral.
O \textbf{currículo} do Ensino Médio será definido pela Base Nacional Comum Curricular ( BNCC ).
Tudo o que será lecionado vai estar dentro do que \textbf{preconizam} os seguintes “ \textbf{itinerários} \textbf{formativos} ” : As escolas da \textbf{rede} não serão obrigadas a oferecer aos alunos todos os cinco \textbf{itinerários}, mas obrigatoriamente um deles para cumprir o que determina a Lei nº 13.415/2017 ( 60 \% da \textbf{carga} \textbf{horária} por conteúdos da BNCC e 40 \% optativos na parte diversificada ), ou seja : Formação Geral Básica ( Base Nacional Comum ) : aumenta -se o número de aulas da matriz comum para algumas disciplinas específicas ( língua portuguesa e matemática ).
Neste também consta o aprofundamento do \textbf{itinerário}, que serão os componentes do \textbf{curso} \textbf{técnico} / FIC ou Programa estudado pelo aluno.
Fica a critério de cada escola decidir por \textbf{ofertar} Ensino Técnico, FIC ou Programa de Aprendizagem. \textbf{Itinerários} \textbf{Formativos} ( Parte diversificada ) : os alunos terão a oportunidade de ter contato com atividades voltadas para diversas áreas de maneira integrada, sem qualquer separação.
O Projeto de Vida será componente obrigatório assim como as competências \textbf{eletivas} ( as quais o aluno escolhe o que quer estudar – optativas, desde que haja total relação com o Projeto de Vida do aluno ).
Sobre esse aspecto, vale destacar a possibilidade de \textbf{oferta} de \textbf{qualificações} e \textbf{especializações} \textbf{técnicas} como partes integrantes dos \textbf{itinerários} \textbf{formativos}.
A organização de \textbf{cursos} \textbf{técnicos} com \textbf{certificações} intermediárias de \textbf{qualificação} \textbf{profissional} \textbf{técnica} consiste em certificar o aluno que \textbf{cursar} com aprovação uma etapa do \textbf{curso} que seja reconhecida como ocupação no mundo do trabalho.
Essa possibilidade se baseia no princípio da flexibilidade e no desenvolvimento e aprimoramento contínuo de competências, que devem orientar a organização dos \textbf{currículos} em diferentes perspectivas : na \textbf{oferta} de \textbf{cursos}, na composição das Unidades Curriculares, nas \textbf{certificações} intermediárias e no aproveitamento de estudos.
Ao incluir o aproveitamento das Unidades Curriculares no âmbito dos \textbf{itinerários} \textbf{formativos}, percebe -se o esforço em criar estratégias intercambiáveis de formação que respondam, de um lado, às \textbf{trajetórias} \textbf{profissionais} dos alunos e, de outro, à complexidade e rapidez do mundo do trabalho.
Assim, \textbf{amparado} nas bases legais que regem a educação \textbf{profissional} brasileira, a SEED/AP define \textbf{itinerário} \textbf{formativo} como um conjunto de percursos de formação propiciados por uma instituição de educação \textbf{profissional}, em cada um dos diferentes eixos tecnológicos.
Com o potencial de promover o aproveitamento contínuo e a articulação entre os eixos tecnológicos e segmentos de mercado, o \textbf{itinerário} \textbf{formativo} torna -se um instrumento importante para o planejamento de vida e \textbf{carreira} do \textbf{trabalhador}, uma vez que amplia suas oportunidades de (re)inserção no mercado, de promoção e de \textbf{mobilidade} \textbf{profissional}.
Ao mesmo tempo, oferece às empresas um mapa para seus investimentos em formação de pessoas e, às instituições formadoras, fornece um roteiro estratégico para a organização da \textbf{oferta} de educação \textbf{profissional}.
É importante reforçar que o \textbf{itinerário} \textbf{formativo} deriva da análise do \textbf{itinerário} \textbf{profissional} de cada área de atuação.
Assim, enquanto o \textbf{itinerário} \textbf{formativo} corresponde ao conjunto de \textbf{cursos} e programas que a Instituição oferece aos alunos para que possam planejar suas \textbf{carreiras}, o \textbf{itinerário} \textbf{profissional} equivale ao conjunto de ocupações com identidades definidas no mundo do trabalho em determinada área \textbf{profissional}.
Para identificá -lo, é preciso considerar o escopo das ocupações, o campo de atuação, a \textbf{legislação}, a interface com outras ocupações e os limites funcionais entre elas.
Dado que cada \textbf{perfil} \textbf{profissional} exige competências com graus específicos de complexidade, o \textbf{itinerário} \textbf{profissional} oferece informações fundamentais para a definição da amplitude do \textbf{itinerário} \textbf{formativo}, permitindo mapear os \textbf{requisitos} e percursos \textbf{formativos} necessários para atuar em determinada ocupação.
Por meio dos \textbf{itinerários} \textbf{formativos} o aluno pode situar seu contexto \textbf{profissional} e definir seu percurso \textbf{formativo}, podendo tanto aumentar o grau de sua formação na mesma área de atuação como fazer \textbf{cursos} em outras áreas complementares ou compatíveis.
Dessa forma, os alunos podem escolher entre diferentes níveis e modalidades de educação \textbf{profissional} disponíveis, de acordo com seus níveis de \textbf{escolaridade} e experiências de vida.
É dentro desse campo de possibilidades de formação e de vivência no mundo do trabalho que se constituem as \textbf{trajetórias} de profissionalização de cada indivíduo.

% matched lemmas: amparar, carga, carreira, catálogo, certificação, currículo, cursar, curso, diretriz, eletivo, escolaridade, especialização, estadual, estruturante, formativo, horário, itinerário, legislação, mobilidade, oferta, ofertar, perfil, preconizar, profissional, qualificação, rede, reforma, requisito, trabalhador, trajetória, transversalidade, técnico, órgão
\end{textsample}
